% ------------------------------------------------------------------------
%
\begin{frame}{Higher-order rewrite systems}

  We already extended term-rewriting systems (TRS) with conditions.

  \bigskip

  We can also extend them in another direction: variables could denote
  functions. These are called \textbf{higher\hyp{}order systems}.

  \bigskip

  They enable defining functions that can be parameterised by other
  functions.

  \bigskip

  For example, let us consider \textbf{iterators}. Typically, those
  are functions that traverse a data structure and apply a given
  function at each step.

\end{frame}

% ------------------------------------------------------------------------
%
\begin{frame}{Maps on stacks}

  In the case of stacks, a simple iterator is a function `\fun{map}'
  that takes a stack and returns a stack made of the images of the
  elements, in the original order. So $[e_1; e_2; \ldots; e_n]$ becomes $[f(e_1);
    f(e_2); \ldots; f(e_n)]$.
  \begin{equation*}
    \begin{array}{r@{\;}c@{\;}l@{}}
      \fun{map}(f, \el) &\smashedrightarrow{\alpha} & \el;\\
      \fun{map}(f, \cons{x}{s}) &\smashedrightarrow{\beta}
      & \cons{f(x)}{\fun{map}(f, s)}.
    \end{array}
  \end{equation*}

  For example, if we set
  \begin{align*}
    \fun{incr}(n)   &\rightarrow n + 1.\\
    \fun{square}(n) &\rightarrow n * n.
  \end{align*}
  then we have
  \begin{align*}
    \fun{map}(\fun{incr}, [4; 4; 5; 0]) &\twoheadrightarrow [5; 5; 6; 1].\\
    \fun{map}(\fun{square}, [4; 4; 5; 0]) &\twoheadrightarrow [16; 16; 25; 0].
  \end{align*}

\end{frame}

% ------------------------------------------------------------------------
%
\begin{frame}{Folds on stacks}

  \textbf{Folds} take an additional parameter with is supposed to be
  the initial value of an \textbf{accumulator}. In this framework, the
  given function is applied to an element of the stack \emph{and to
  the current value of the accumulator}.

  \bigskip

  For example, for summing all integers in a stack, we would have
  \begin{equation*}
    \fun{fold} (\fun{add}, [4; 4; 5; 0], 0) \twoheadrightarrow 13,
  \end{equation*}
  where $\fun{add}(x,y) \rightarrow x + y$.

  \bigskip

  Note how the direction of iteration does not matter, because
  \begin{equation*}
    4 + (4 + (5 + (0 + \underline{0})))
    = (((\underline{0} + 4) + 4) + 5) + 0,
  \end{equation*}
  where the initial value of the accumulator is~\(\underline{0}\).

\end{frame}

% ------------------------------------------------------------------------
%
\begin{frame}{Lambdas}

  Note how cumbersome it is sometimes to have to define the function
  applied by the iterator (here a fold), when it is probably used only
  once.

  \bigskip

  To remedy this, we can further extend or systems with
  \textbf{anonymous functions}, also known as \textbf{lambdas} (a
  name coming from the $\lambda$\hyp{}calculus, which is a different
  framework in which to model computations.)

  \bigskip

  For example, we will write $\Xfun\; x \rightarrow x$ to denote the
  identity function, without naming it. Of course, we have to forbid
  calling any function `\fun{fun}` now (in other words, `\Xfun' is a
  \textbf{keyword}).

  \bigskip

  In our case above:
  \begin{equation*}
    \fun{fold} ((\Xfun\; (x,y) \rightarrow x + y), [4; 4; 5; 0], 0)
    \twoheadrightarrow 13,
  \end{equation*}

\end{frame}

% ------------------------------------------------------------------------
%
\begin{frame}{A map is a fold}

  A map is then a special case of a fold, where the accumulator is the
  new stack:
  \begin{equation*}
    \fun{map}(f, s)
    \rightarrow \fun{fold}((\Xfun\; (e, a) \rightarrow \cons{f(e)}{a}), s, \el),
  \end{equation*}
  where $a$ is the accumulator and $e$ a stack element.

  \bigskip

  Contrary to $(+)$, pushing on a stack is not a symmetric relation,
  so the order of traversal matters.

  \bigskip

  Here, since we want to conserve the order of the original elements
  of the stack, the fold must apply the given function first to the
  bottommost element, then the penultimate, then the antepenultimate
  etc. until the top. If we consider a stack a being layout
  horizontally, this means applying the function from right to left.

  \bigskip

  This is called a \textbf{right fold}.

\end{frame}

% ------------------------------------------------------------------------
%
\begin{frame}{Right fold}

  Let's rename `\fun{fold}' as `\fun{foldr}' (`fold right'). Then:
  \begin{equation*}
    \begin{array}{r@{\;}c@{\;}l@{}}
      \fun{foldr} (f, [], a) & \rightarrow & a;\\
      \fun{foldr} (f, \cons{x}{s}, a) & \rightarrow & f(x, \fun{foldr}(f, s, a)).
    \end{array}
  \end{equation*}
  A way to understand this definition is by extension:
  \begin{equation*}
    \fun{foldr}(f,[e_1; e_2; \ldots; e_n], a) \equiv f(e_1, f(e_2, \ldots f(e_n, a))).
  \end{equation*}

  Note that the recursive call in the definition of `\fun{foldr}' is
  not in tail position. (See page~\pageref{tail_call} an onwards.)

  \bigskip

  In this case, this means that all the calls to the iterated
  function~$f$ will be waiting for all the calls to~`\fun{foldr}' to
  be finished (as well as all the calls to~$f$ on the previous
  elements of the stack).

  \bigskip

  Generally speaking, the call stack will grow at least proportionally
  to the length of the stack.

\end{frame}

% ------------------------------------------------------------------------
%
\begin{frame}{Left fold}

  The dual iteration of a right fold is a \textbf{left fold}:
  \begin{equation*}
    \fun{foldr}(f,[e_1; e_2; \ldots; e_n], a) \equiv f(e_n, \ldots f(e_2, f(e_1, a))).
  \end{equation*}
  It is a simple matter of rearranging the calls in the definition of `\fun{foldr}`:
  \begin{equation*}
    \begin{array}{r@{\;}c@{\;}l@{}}
      \fun{foldl} (f, [], a) & \rightarrow & a;\\
      \fun{foldl} (f, \cons{x}{s}, a) & \rightarrow &\fun{foldl}(f, s, f(x,a)).
    \end{array}
  \end{equation*}

  Note how the recursive call is in tail position, so, if tail call
  optimisation is performed, the call stack will grow a constant a
  constant amount if the calls to the given function also grow the
  stack a constant amount (in total).

  \bigskip

  These considerations on memory usage are only meaningful when
  expressing our systems in programming languages and running the
  programs, but not when designing a system.

\end{frame}
